% Part: incompleteness
% Chapter: incompleteness-provability
% Section: second-incompleteness-thm

\documentclass[../../../include/open-logic-section]{subfiles}

\begin{document}

\olfileid{inc}{inp}{2in}

\olsection{د ناتکميلۍ دويمه قضيه}

دا دعوه څنګه بيان کړو چې~$\Th{PA}$ خپله سازګاري نۀ شي ثابتولے؟
دا ويل چې~$\Th{PA}$ ناسازګاره ده، د دې ويلو برابر دي چې
$\Th{PA} \Proves \eq[0][1]$۔ نو د سازګارۍ جمله~$\OCon[\Th{PA}]$
د !!{sentence}~$\lnot
\OProv[\Th{PA}](\gn{\eq[0][1]})$ په توګه اخيستے شو، او بيا
لاندې قضيه اړينه پايله راکوي:

\begin{thm}
\ollabel{thm:second-incompleteness}
که~$\Th{PA}$ سازګاره وي، نو~$\Th{PA}$،
$\OCon[\Th{PA}]$ نۀ !!{derive} کوي۔
\end{thm}

دا خبره مهمه ده چې قضيه د~$\OCon[\Th{PA}]$ په ټاکلي تمثيل پورې
تړلې ده، يعنې د~$\OProv[\Th{PA}](y)$ په ټاکلي تمثيل پورې۔ موږ
يوازې دومره کاروو چې د~$\OProv[\Th{PA}](y)$ تمثيل د
!!{derivability} درې شرطونه پوره کوي؛ نو قضيه هرې هغې تيورۍ ته
غځېږي چې د !!a{derivability} پريديکات ئې دغه خاصيتونه ولري۔

د ګوډل د استدلال لنډه نقشه لوستل ګټور دي، ځکه قضيه پکې لکه د ښې
ټوکې اغېزمنه وروستۍ جمله راځي۔ استدلال داسې دے: $!G_\Th{PA}$
دې د ګوډل هغه جمله وي چې د~\olref[1in]{thm:first-incompleteness}
په ثبوت کښې مو جوړه کړه۔ موږ وښودله: «که~$\Th{PA}$ سازګاره وي،
نو~$\Th{PA}$، $!G_\Th{PA}$ نۀ !!{derive} کوي»۔ که دا خبره
پخپله \emph{په}~$\Th{PA}$ کښې صوري کړو، لاندې ثبوت لرو:
\[
\OCon[\Th{PA}] \lif \lnot \OProv[\Th{PA}](\gn{!G_\Th{PA}}).
\]
اوس فرض کړئ~$\Th{PA}$، $\OCon[\Th{PA}]$ !!{derive}s کوي۔ بيا
$\lnot
\OProv[\Th{PA}](\gn{!G_\Th{PA}})$ هم !!{derive}s کوي۔
\textbf{د ژباړې د نسخې سمون (OLPRV-002):} اصلي سرچينه په دې
جمله کښې~$\Prov[\Th{PA}]$ ليکي، خو د صوري جملې دپاره
$\OProv[\Th{PA}]$ پکار دے؛ همدا نښه پورته په ثبوتېدونکې
پايلي کښې راغلې ده۔ ځکه چې~$!G_\Th{PA}$ د ګوډل جمله ده، دا له
$!G_\Th{PA}$ سره معادله ده۔ نو~$\Th{PA}$، $!G_\Th{PA}$
!!{derive}s کوي۔

خو موږ پوهېږو چې که~$\Th{PA}$ سازګاره وي، $!G_\Th{PA}$ نۀ
!!{derive} کوي!{} نو که~$\Th{PA}$ سازګاره وي،
$\OCon[\Th{PA}]$ هم نۀ !!{derive} کولے شي۔

د استدلال د لا کره کولو دپاره~$!G_\Th{PA}$ د~$\Th{PA}$ د ګوډل
جمله اخلو او د !!{derivability} شرطونه (P1)--(P3) کاروو، څو
وښيو چې~$\Th{PA}$، $\OCon[\Th{PA}] \lif
!G_\Th{PA}$ !!{derive}s کوي۔ له دې به څرګنده شي چې~$\Th{PA}$،
$\OCon[\Th{PA}]$ نۀ !!{derive} کوي۔ دا په~$\Th{PA}$ کښې د
ثبوت لنډه نقشه ده۔ (د آسانتيا دپاره د~$\Th{PA}$ لاندې ليکونه
پرېږدو۔)
\begin{align}
& !G \liff \lnot \OProv(\gn{!G}) \ollabel{G2-1}\\
& \qquad\text{$!G$ د ګوډل جمله ده}\notag \\
& !G \lif \lnot \OProv(\gn{!G}) \ollabel{G2-2}\\
  & \qquad\text{له \olref{G2-1} څخه} \notag\\
& !G \lif
  (\OProv(\gn{!G}) \lif \lfalse) \ollabel{G2-3}\\
  & \qquad\text{له \olref{G2-2} څخه د منطق له مخې}\notag\\
& \OProv(\gn{
    !G \lif
    (\OProv(\gn{!G}) \lif \lfalse)
  }) \ollabel{G2-4}\\
  & \qquad\text{له \olref{G2-3} څخه د P1 شرط له مخې} \notag\\
& \OProv(\gn{!G}) \lif
  \OProv(\gn{
    (\OProv(\gn{!G}) \lif \lfalse)
    }) \ollabel{G2-5}\\
  & \qquad\text{له \olref{G2-4} څخه د P2 شرط له مخې} \notag\\
& \OProv(\gn{!G}) \lif (\OProv(\gn{\OProv(\gn{!G})}) \lif \OProv(\gn{\lfalse})) \ollabel{G2-6}\\
  & \qquad\text{له \olref{G2-5} څخه د P2 شرط او منطق له مخې} \notag\\
& \OProv(\gn{!G}) \lif
  \OProv(\gn{\OProv(\gn{!G})}) \ollabel{G2-7}\\
   & \qquad\text{د P3 له مخې} \notag\\
& \OProv(\gn{!G}) \lif \OProv(\gn{\lfalse}) \ollabel{G2-8}\\
  & \qquad \text{له \olref{G2-6} او \olref{G2-7} څخه د منطق له مخې}\notag\\
& \OCon \lif \lnot \OProv(\gn{!G}) \ollabel{G2-9}\\
  & \qquad\text{د \olref{G2-8} متقابل عکس او $\OCon \ident \lnot \OProv(\gn{\lfalse})$}\notag \\
& \OCon \lif !G \notag\\
  & \qquad\text{له \olref{G2-1} او \olref{G2-9} څخه د منطق له مخې}\notag
\end{align}
په پورته ثبوت کښې د منطق کارول يوازې د قضيې د منطق بنسټيز
حقايق دي: مثلاً~\olref{G2-3} له
$\Proves \lnot!A \liff (!A\lif
\lfalse)$ څخه کار اخلي، او~\olref{G2-8} له
$!A \lif (!B \lif !C), !A \lif !B
\Proves !A \lif !C$ څخه۔ په~\olref{G2-5} او~\olref{G2-6}
کښې د P2 شرط کارول د هغه پر بېلګو
$\OProv(\gn{!A \lif !B}) \lif
(\OProv(\gn{!A}) \lif \OProv(\gn{!B}))$ ولاړ دي۔ په لومړۍ
بېلګه کښې $!A \ident
!G$ او $!B \ident \OProv(\gn{!G}) \lif \lfalse$ دي؛ په
دويمه کښې $!A
\ident \OProv(\gn{!G})$ او $!B \ident \lfalse$ دي۔
\textbf{د ژباړې د نسخې سمون (OLPRV-003):} د دويمې بېلګې د
$!A$ په اصلي نښه کښې~$\gn{G}$ ليکل شوي وو؛ دلته د لومړۍ
بېلګې او د ثبوت د فورمولونو په شان د جملې نښه~$!G$ ده، نو
$\gn{!G}$ ساتل پکار دي۔

د ناتکميلۍ د دويمې قضيې لا انتزاعي بڼه دا ده:

\begin{thm}
\ollabel{thm:second-incompleteness-gen} فرض کړئ~$\Th{T}$ هره
سازګاره، !!{axiomatized} تيوري ده چې~$\Th{Q}$ غځوي، او
$\OProv[\Th{T}](y)$ هر هغه فورمول دے چې د~$\Th{T}$ دپاره د
!!{derivability} شرطونه P1--P3 پوره کوي۔ نو~$\Th{T}$،
$\OCon[T]$ نۀ !!{derive} کوي۔
\end{thm}

\begin{prob}
وښيئ چې~$\Th{PA}$، $!G_{\Th{PA}} \lif \OCon[\Th{PA}]$
!!{derive}s کوي۔
\end{prob}

\begin{digress}
د دې کيسې درس دا دے چې د رياضي هېڅ «معقوله» سازګاره تيوري
خپله د سازګارۍ جمله نۀ شي !!{derive} کولے۔ فرض کړئ~$\Th{T}$
د رياضي داسې تيوري ده چې~$\Th{Q}$ او د هېلبرټ «متناهي‌پاله»
استدلال (هر څه چې مطلب ترې وي) رانغاړي۔ نو ټوله~$\Th{T}$ د
$\Th{T}$ د سازګارۍ جمله نۀ شي !!{derive} کولے؛ ځکه ئې متناهي‌پاله
برخه هم په لا پياوړي دليل سره د~$\Th{T}$ د سازګارۍ جمله نۀ شي
!!{derive} کولے۔ په دې معنا د «ټولې رياضي» دپاره د سازګارۍ
متناهي‌پاله ثبوت نشته۔

د «متناهي‌پاله» لفظ په تفسير کښې لږ ګنجايش شته۔ ګوډل په خپلې
۱۹۳۱ز مقالې کښې دا شونتيا مني چې هغه څه چې موږ ئې «متناهي‌پاله»
بولو، ښايي د رياضي له هغو ډولونو بهر وي چې هېلبرټ غوښتل صوري
کړي۔ خو ګوډل په دې کښې د نرمۍ څخه کار اخيست؛ نن دا ګرانه ده چې
داسې څه ومومو چې په معقول ډول متناهي‌پاله وبلل شي، خو مثلاً
په~$\Th{ZFC}$، يعنې د ټاکنې له بديهي اصل سره د زرملو--فرانکل
د سټونو په تيورۍ کښې، صوري نۀ شي۔
\end{digress}

\end{document}
